Now that we have some credentials, we can use them to access the server.
This is done with SSH, which listens on port 2022.\\
\pandocbounded{\includegraphics[keepaspectratio,alt={login}]{images/login.png}}

Now that we have signed in, we need to perform reconnaissance. The entry
message says that we have ``a few elevated maintenance rights'' so we
should check what we have access to. This involves checking the
permissions of files in our home directory (with \texttt{ls\ -al}),
looking for any setuid binaries (with
\texttt{find\ /\ -type\ f\ -perm\ -4000\ 2\textgreater{}/dev/null}),
looking for any capabilities (with
\texttt{getcap\ -r\ /\ 2\textgreater{}/dev/null}), and looking for any
binaries we can run as root (with \texttt{sudo\ -l}).\\
\pandocbounded{\includegraphics[keepaspectratio,alt={privileges}]{images/privileges.png}}

There isn't anything interesting in our home directory, but it looks
like there are a few binaries with the setuid bit enabled and a single
binary (\texttt{find}) that can be run using \texttt{sudo} with no
password. This is actually the most interesting one, and is worth
exploring further. We can look at what command line arguments it accepts
using \texttt{find\ -\/-help}, as it may have something that can be
exploited.\\
\pandocbounded{\includegraphics[keepaspectratio,alt={help}]{images/help.png}}

It looks like it supports an \texttt{-exec} flag, which can seemingly
run arbitrary commands. We can test this out by crafting a payload that
causes the command to run \texttt{/bin/bash}, which should give us a
root shell when running with \texttt{sudo} (as \texttt{sudo} makes it
run as root and no password is required to run \texttt{find} using
\texttt{sudo}). This only needs to be a simple payload that follows this
command's conventions, so
\texttt{sudo\ find\ -exec\ /bin/bash\ \textbackslash{};} should
hopefully work.\\
\pandocbounded{\includegraphics[keepaspectratio,alt={exec}]{images/exec.png}}

This has given us a root shell, meaning this attack has been completed
successfully.
